% =====================================================================
\chapter{Cuantificar el valor y el riesgo}
\label{tema:5}
% =====================================================================

Un optimizador se justifica si su valor esperado supera a su coste y si sus riesgos son aceptables y están controlados. Este tema da las herramientas para afirmarlo con números: intervalos de confianza para la mejora con diseños pareados (capítulo~2 de Montgomery~\cite{montgomery}) y con \emph{bootstrap} (lectura~24 del curso 18.05 del MIT~\cite{mit1805}), tamaño de un piloto, valor actual neto (sesiones «Present Value Relations» y «Capital Budgeting» del curso 15.401 del MIT~\cite{mit15401}), medidas de riesgo de cola (Rockafellar y Uryasev~\cite{ru}), análisis de modos de fallo (IEC~60812~\cite{iec60812}) y el marco del Reglamento europeo de inteligencia artificial~\cite{aiact,omnibus}. Se apoya en el módulo~0.9 (temas~4 y~5), en el contraste del ahorro del módulo~26 (tema~6, ejemplo «Contraste del ahorro») y en los experimentos en línea de Kohavi, Tang y Xu~\cite{kohavi}.

\section{Coste y beneficio}

\begin{definicion}[Valor actual neto]\label{def:van}
Un proyecto con inversión inicial $I$ y flujos netos $B_t-O_t$ (beneficios menos costes de operación) al final de cada año $t=1,\dots,N$, con tipo de descuento anual $\rho>0$, tiene \emph{valor actual neto}
\[
\VAN=-I+\sum_{t=1}^{N}\frac{B_t-O_t}{(1+\rho)^{t}}.
\]
La \emph{tasa interna de rentabilidad} (TIR) es el $\rho$ que anula el VAN, y el \emph{plazo de recuperación}, el primer año en que los flujos acumulados sin descontar cubren $I$.
\end{definicion}

\begin{proposicion}[Anualidad]\label{prop:anualidad}
Si $B_t-O_t=b$ es constante, $\VAN=-I+b\,a_{N,\rho}$ con $a_{N,\rho}=\sum_{t=1}^{N}(1+\rho)^{-t}=\dfrac{1-(1+\rho)^{-N}}{\rho}$.
\end{proposicion}

Los beneficios de un optimizador se miden \emph{respecto a una alternativa}, y la elección de la alternativa decide el resultado. La alternativa honrada no es la práctica actual, sino la mejor solución sencilla que se adoptaría sin IA. En el mantenimiento del \cref{tema:4}, pasar de intervenir cada 70 días a hacerlo cada 30 (política~B) no requiere sensores ni modelos de clasificación: basta un análisis de fiabilidad. El valor atribuible al optimizador con IA es la diferencia entre la política~D y la~B.

\section{Intervalos de confianza para la mejora}

\begin{proposicion}[Ventaja del diseño pareado]\label{prop:pareado}
Sean $Y_A$ e $Y_B$ los costes de dos políticas en el mismo periodo, con varianzas $\sigma_A^{2}$, $\sigma_B^{2}$ y correlación $r$. La diferencia pareada $D=Y_A-Y_B$ tiene varianza $\sigma_A^{2}+\sigma_B^{2}-2r\sigma_A\sigma_B$, menor que la de la diferencia de dos medidas independientes, $\sigma_A^{2}+\sigma_B^{2}$, si y solo si $r>0$.
\end{proposicion}

Comparar dos políticas en los mismos días, o en el gemelo con los mismos números aleatorios (módulo~26, tema~5, §5.1), crea esa correlación positiva: la carga, el desgaste y el azar comunes se cancelan. Con $n$ diferencias pareadas de media $\bar d$ y desviación típica $s_d$, el intervalo de confianza de nivel $1-\alpha$ para la mejora media es $\bar d\pm t_{n-1,1-\alpha/2}\,s_d/\sqrt n$ (módulo~0.9, tema~4, §4.4).

\begin{definicion}[Intervalos \emph{bootstrap}]\label{def:bootstrap}
Sean $d_1,\dots,d_n$ las diferencias observadas y $\bar d$ su media. Se generan $B$ remuestras de tamaño $n$ con reemplazamiento, con medias $\bar d^{*}_1,\dots,\bar d^{*}_B$. El intervalo \emph{percentil} de nivel $1-\alpha$ es $[\bar d^{*}_{(\alpha/2)},\bar d^{*}_{(1-\alpha/2)}]$, con los cuantiles empíricos de las medias remuestreadas. El intervalo \emph{básico} usa $\delta^{*}=\bar d^{*}-\bar d$: $[\bar d-\delta^{*}_{(1-\alpha/2)},\ \bar d-\delta^{*}_{(\alpha/2)}]$~\cite{mit1805}.
\end{definicion}

El \emph{bootstrap} no exige normalidad, lo que importa aquí: el coste semanal de una línea es casi siempre pequeño y, de vez en cuando, salta en 9000~€ por una avería. Ambos métodos suponen observaciones independientes; si las semanas están autocorrelacionadas hay que remuestrear por bloques.

\begin{proposicion}[Tamaño de un piloto]\label{prop:tamano}
Para detectar una mejora media $\delta>0$ con un contraste bilateral de nivel $\alpha$ y potencia $1-\beta$ con diferencias pareadas de desviación típica $\sigma_D$, hacen falta aproximadamente
\[
n\approx\Bigl(\frac{(z_{1-\alpha/2}+z_{1-\beta})\,\sigma_D}{\delta}\Bigr)^{2}
\]
pares (módulo~0.9, tema~5, §5.4).
\end{proposicion}

\section{Caso de estudio: ¿cuánto vale el optimizador de mantenimiento?}

\addtocounter{definicion}{1}%

\section{Robustez y riesgo de cola}

La media no lo es todo: a una planta le preocupa el año malo. Para una pérdida $L$ con distribución continua, el \emph{valor en riesgo} al nivel $\alpha$ es el cuantil $\mathrm{VaR}_\alpha(L)=F_L^{-1}(\alpha)$, y el \emph{valor en riesgo condicional} es la media de la cola, $\mathrm{CVaR}_\alpha(L)=\E[L\mid L\ge\mathrm{VaR}_\alpha(L)]$.

\begin{proposicion}[Fórmula de Rockafellar y Uryasev]\label{prop:cvar}
Si $L$ tiene densidad y esperanza finita,
\[
\mathrm{CVaR}_\alpha(L)=\min_{\tau\in\R}\Bigl\{\tau+\frac{1}{1-\alpha}\E\bigl[(L-\tau)^{+}\bigr]\Bigr\},
\]
y el mínimo se alcanza en $\tau=\mathrm{VaR}_\alpha(L)$.
\end{proposicion}

La fórmula convierte la minimización del CVaR de una decisión en un programa lineal cuando $L$ se representa con escenarios: basta añadir $\tau$ y una variable $u_s\ge L_s-\tau$, $u_s\ge0$ por escenario~\cite{ru}. En los 20 años simulados del caso, el coste anual de la política~D tiene media $159\,300$~€ y CVaR al 90\,\% (media de los dos peores años) de $186\,300$~€; la política~B, $193\,670$ y $216\,200$~€. D es mejor también en la cola. Las \emph{pruebas de estrés} completan el análisis: repetir la evaluación con un coste de avería un 50\,\% mayor, con un técnico de baja, con un sensor que se congela en su último valor o con una semana de datos perdidos.

\section{Seguridad, supervisión y modos de fallo}

Los principios del módulo~26 (tema~6) se aplican a cualquier optimizador. (i) Las \emph{restricciones duras} (capacidad de los técnicos, inspecciones obligatorias por ley, bloqueos de seguridad de las máquinas) están en el modelo de decisión y se comprueban otra vez, de forma independiente, antes de emitir una orden. (ii) Hay una \emph{política de respaldo} sencilla (aquí, la política por edad) a la que se vuelve si el modelo o los datos fallan. (iii) El despliegue recorre los modos asesor, supervisado y autónomo con supervisión (módulo~26, tema~6, §6.2), y en cada uno se registra quién decidió qué.

\begin{definicion}[Análisis modal de fallos y efectos]\label{def:amfe}
El AMFE (FMEA en inglés; IEC~60812~\cite{iec60812}) enumera los modos de fallo de cada componente, sus efectos y sus causas, y los prioriza. En la variante clásica se puntúan de 1 a 10 la \emph{gravedad} $S$, la \emph{ocurrencia} $O$ y la \emph{detectabilidad} $Dt$ (10 = indetectable), y se ordena por el \emph{número de prioridad de riesgo} $\mathrm{NPR}=S\cdot O\cdot Dt$; la norma admite otras formas de calcularlo y, en todo caso, los modos de gravedad máxima se tratan aunque su NPR sea bajo.
\end{definicion}

\begin{table}[h]
\centering\small
\caption{AMFE resumido del optimizador de mantenimiento.}\label{tab:amfe}
\setlength{\tabcolsep}{4pt}\footnotesize
\begin{tabularx}{\textwidth}{@{}>{\raggedright\arraybackslash}p{2.8cm}>{\raggedright\arraybackslash}X cccc >{\raggedright\arraybackslash}p{3.0cm}@{}}
\toprule
Modo de fallo & Efecto & $S$ & $O$ & $Dt$ & NPR & Medida \\
\midrule
sensor congelado en su último valor & riesgo infravalorado; avería no prevista & 8 & 3 & 7 & 168 & regla de varianza nula; respaldo por edad \\
deriva tras cambio de producto & probabilidades mal calibradas & 6 & 5 & 4 & 120 & PSI semanal; recalibración \\
resolvedor sin solución a tiempo & sin plan el lunes & 5 & 2 & 2 & 20 & plan de la semana anterior; límite de tiempo \\
confianza excesiva del técnico & se ejecutan planes erróneos & 7 & 4 & 6 & 168 & mostrar motivos; auditoría de anulaciones \\
etiquetas mal registradas en el GMAO & modelo entrenado con ruido & 6 & 4 & 6 & 144 & validación de registros; revisión mensual \\
\bottomrule
\end{tabularx}
\end{table}

\section{El Reglamento europeo de inteligencia artificial}

El Reglamento (UE) 2024/1689~\cite{aiact} entró en vigor el 1 de agosto de 2024 y sigue un enfoque basado en el riesgo: prohíbe algunas prácticas, impone requisitos a los sistemas de \emph{alto riesgo}, establece obligaciones de transparencia para otros y deja el resto sin obligaciones específicas. Lo que sigue es una descripción general de su texto público; la clasificación de un sistema concreto es una cuestión jurídica.

\begin{enumerate}[label=\arabic*.]
  \item \emph{Alto riesgo} (artículo~6). Lo es un sistema que sea componente de seguridad de un producto, o el propio producto, regulado por la legislación de armonización del anexo~I y sujeto a evaluación de la conformidad por terceros; y lo son los sistemas del anexo~III, entre ellos los componentes de seguridad en la gestión y operación de infraestructuras críticas, como el suministro de agua, gas, calefacción o electricidad (anexo~III, punto~2). Un sistema del anexo~III no es de alto riesgo si no influye de forma sustancial en la decisión, por ejemplo si realiza una tarea procedimental limitada o preparatoria (artículo~6, apartado~3), salvo que elabore perfiles de personas.
  \item \emph{Requisitos} (artículos~8 a~15): sistema de gestión de riesgos, gobernanza de los datos, documentación técnica, registro automático de eventos, transparencia hacia quien lo despliega, supervisión humana y precisión, robustez y ciberseguridad. El artículo~14 pide que las personas que supervisan puedan entender las capacidades y los límites del sistema y detectar anomalías, sean conscientes del sesgo de automatización, interpreten correctamente sus salidas, puedan decidir no usarlo o anular sus resultados, e intervenir o detenerlo de forma segura.
  \item \emph{Calendario.} El reglamento se aplica con carácter general desde el 2 de agosto de 2026. La modificación conocida como «ómnibus digital», publicada en el Diario Oficial el 24 de julio de 2026 y en vigor desde el 27 de julio, aplaza las obligaciones de alto riesgo: al 2 de diciembre de 2027 para los sistemas del anexo~III y al 2 de agosto de 2028 para los del anexo~I; según la información publicada, la maquinaria pasa a un régimen sectorial propio~\cite{omnibus,aiact}.
\end{enumerate}

Aplicado al caso: un optimizador del calendario de mantenimiento de una línea de envasado, cuyas órdenes revisa un técnico y que no actúa sobre las funciones de seguridad de las máquinas, no parece encajar en ninguna categoría de alto riesgo. El mismo optimizador aplicado a las bombas de una red de abastecimiento de agua, si actuara como componente de seguridad de su operación, sí podría serlo. En ambos casos, la arquitectura de los temas~4 y~5 (registro de decisiones, supervisión humana con capacidad de anular, política de respaldo, monitorización, documentación) cumple de forma natural buena parte de los requisitos, y es buena práctica adoptarla aunque no sea obligatoria. Si se usan datos de personas (por ejemplo, qué técnico anuló qué orden), se aplica además la normativa de protección de datos.

\begin{table}[h]
\centering\small
\caption{Tecnologías del tema~5.}\label{tab:tec5}
\begin{tabularx}{\textwidth}{@{}>{\raggedright\arraybackslash}p{2.5cm}>{\raggedright\arraybackslash}X>{\raggedright\arraybackslash}p{2.1cm}>{\raggedright\arraybackslash}X@{}}
\toprule
Herramienta & Para qué sirve & Licencia o tipo & Cuándo elegirla \\
\midrule
SciPy (\texttt{stats}) & contrastes $t$, intervalos y \emph{bootstrap} (función \texttt{bootstrap}) & abierta (BSD-3) & análisis de pilotos y de simulaciones \\
statsmodels & modelos con errores autocorrelacionados, potencia y tamaño muestral & abierta (BSD-3) & series semanales, diseños con covariables \\
numpy-financial & VAN, TIR y anualidades & abierta (BSD-3) & valoración económica reproducible \\
MLflow~\cite{mlflow} & registrar cada evaluación con su modelo, datos y resultado & abierta (Apache~2.0) & trazabilidad exigida por la documentación técnica \\
Plantillas de AMFE (hoja de cálculo) & análisis de modos de fallo & --- & revisión de riesgos con operación, mantenimiento y seguridad \\
\bottomrule
\end{tabularx}
\end{table}

\addtocounter{definicion}{1}%

